\section{Exact moments at a finite update count}\label{sec:discrete}

This section uses the population response model
$y=x^\top\beta^*+\varepsilon$ in a fixed dimension, with known population
centering $\E x=0$ and fixed positive diagonal scales $S$. There is no
intercept update. Assume
\[
 C=\E xx^\top\succ0,\qquad \E\|x\|^4<\infty,\qquad
 \E(\varepsilon\mid x)=0,\qquad
 \E\varepsilon^2<\infty,\qquad
 \E[\varepsilon^2\|x\|^2]<\infty.
\]
The last condition controls gradient noise and does not follow from finite
noise variance alone. It permits heteroscedastic noise and does not require
Gaussian errors or independence of $x$ and $\varepsilon$.
Write $P=S^{-2}$, and use the scaled random feature $z=S^{-1}x$ and
scaled coefficient error
\[
 H=\E zz^\top=S^{-1}CS^{-1},\qquad
 u_k=S(\beta_k-\beta^*),\qquad R=I-\eta H.
\]
Here $z$ denotes a random feature vector, rather than the fixed empirical
score in the setup. Let the initialization have finite second moment and be
independent of all training observations. Deterministic initialization,
including zero, is a special case.

\subsection{Population gradient descent and the online mean}

Population gradient descent on the averaged squared loss obeys
\[
 \beta_{k+1}=\beta_k-\eta PC(\beta_k-\beta^*),\qquad u_{k+1}=Ru_k.
\]
Consequently, for deterministic $\beta_0$ and every integer $k\ge0$,
\begin{equation}\label{eq:discrete-gd}
 \beta_k=\beta^*+(I-\eta PC)^k(\beta_0-\beta^*).
\end{equation}
For random initialization with scaled mean error $m_0$, covariance
$\Sigma_0$, and second moment $U_0$, its mean error is $R^km_0$, its covariance
is $R^k\Sigma_0R^k$, and its second moment is $R^kU_0R^k$.
Thus both matrix moments obey $V_{k+1}=RV_kR$ in population GD.
There is no observation-noise contribution in this
population update. These identities hold for every $\eta>0$ at finite $k$.
Convergence for every initial value holds precisely when
$0<\eta<2/\lambda_{\max}(H)$, since $H$ is symmetric positive definite.

For fresh-batch SGD, let $(x_{ki},\varepsilon_{ki})$, $1\le i\le B$, be
independent copies of $(x,\varepsilon)$, independently across updates and
independently of the initialization. Define
\[
 \widehat H_k=\frac1B\sum_{i=1}^B z_{ki}z_{ki}^\top,\qquad
 \xi_k=\frac1B\sum_{i=1}^B z_{ki}\varepsilon_{ki}.
\]
The averaged stochastic gradient gives
\begin{equation}\label{eq:discrete-sgd-update}
 u_{k+1}=(I-\eta\widehat H_k)u_k+\eta\xi_k,
 \qquad m_k:=\E u_k=R^km_0.
\end{equation}
Thus SGD has exactly the population GD mean at every finite update count.
This identity uses fresh sampling and fixed centering and scaling; it does
not describe repeated passes through a fixed noisy sample. Comparing batch
sizes at a fixed observation budget $N$ requires $k=\lfloor N/B\rfloor$
in these formulas, not a common update count.

\subsection{Second moments and centered covariance}

For symmetric matrices $U$, define
\begin{align}
 \mathcal M(U)&=\E[zz^\top Uzz^\top],\qquad
 \mathcal V(U)=\mathcal M(U)-HUH,
 \qquad Q=\E[\varepsilon^2zz^\top],\label{eq:discrete-operators}\\
 \mathcal T_B(U)&=RUR+\frac{\eta^2}{B}\mathcal V(U)
 =\E[(I-\eta\widehat H_k)U(I-\eta\widehat H_k)].\label{eq:discrete-T}
\end{align}
The identity
$\mathcal V(U)=\E[(zz^\top-H)U(zz^\top-H)]$ shows that $\mathcal V$
maps positive semidefinite matrices to positive semidefinite matrices.
The same is true of $\mathcal T_B$. The fourth moments enter through
$\mathcal M$, including mixed moments between common and rare coordinates.

\begin{proposition}[Finite-update error moments]\label{prop:discrete-moments}
Under the assumptions above, put
$U_k=\E[u_ku_k^\top]$ and
$\Sigma_k=\Cov(u_k)=U_k-m_km_k^\top$.
For every integer $k\ge0$,
\begin{align}
 U_{k+1}&=\mathcal T_B(U_k)+\frac{\eta^2}{B}Q,
 &U_k&=\mathcal T_B^k(U_0)
       +\frac{\eta^2}{B}\sum_{j=0}^{k-1}\mathcal T_B^j(Q),
       \label{eq:discrete-second}\\
 \Sigma_{k+1}&=\mathcal T_B(\Sigma_k)
       +\frac{\eta^2}{B}\{\mathcal V(m_km_k^\top)+Q\}.
       \label{eq:discrete-covariance}
\end{align}
In particular, the centered covariance has the finite-sum decomposition
\begin{equation}\label{eq:discrete-covariance-sum}
 \Sigma_k=\mathcal T_B^k(\Sigma_0)
 +\frac{\eta^2}{B}\sum_{j=0}^{k-1}
    \mathcal T_B^{k-1-j}\{\mathcal V(R^jm_0m_0^\top R^j)\}
 +\frac{\eta^2}{B}\sum_{j=0}^{k-1}\mathcal T_B^j(Q).
\end{equation}
Empty sums are zero, and the zeroth operator power is the identity.
Original-scale quantities are
\begin{equation}\label{eq:discrete-original-moments}
 \E\beta_k-\beta^*=S^{-1}m_k,\qquad
 \Cov(\beta_k)=S^{-1}\Sigma_kS^{-1},\qquad
 \E[(\beta_k-\beta^*)(\beta_k-\beta^*)^\top]=S^{-1}U_kS^{-1}.
\end{equation}
\end{proposition}
\begin{proof}
Conditional on all features in a batch, $\E(\xi_k\mid z_{k1},\ldots,z_{kB})=0$.
Independence of the sampled pairs and the conditional noise-mean assumption
therefore eliminate both cross terms in the second moment of
\eqref{eq:discrete-sgd-update}, even though $\xi_k$ and $\widehat H_k$
need not be independent. They also give
$\E\xi_k\xi_k^\top=Q/B$.
Independence within a batch yields
\[
 \E[\widehat H_kU\widehat H_k]
   =\left(1-\frac1B\right)HUH+\frac1B\mathcal M(U).
\]
Expanding the two update factors proves \eqref{eq:discrete-T}.
Independence from the past then proves the second-moment recurrence.
Subtract $m_{k+1}m_{k+1}^\top=Rm_km_k^\top R$ to obtain
\eqref{eq:discrete-covariance}. Iteration gives both finite sums. The stated
moment assumptions and independence ensure finite moments at every fixed
$k$ by induction. Conjugation by $S^{-1}$ proves the original-scale formulas.
\end{proof}

Equation~\eqref{eq:discrete-covariance-sum} separates propagated initial
uncertainty, multiplicative update variance caused by mean error, and
observation noise. The operator $\mathcal T_B$ also propagates existing
variance through random updates, so the final noise contribution includes
that amplification. With deterministic initialization, $\Sigma_0=0$;
nevertheless $U_0=m_0m_0^\top$ usually is nonzero. In the noiseless case
$Q=0$, SGD may still have positive covariance, whereas initialization at
$\beta^*$ gives no subsequent error. With homoscedastic conditional noise
variance $\sigma_\varepsilon^2$, $Q=\sigma_\varepsilon^2H$.
For an independent test feature from the training population, expected
excess squared loss is
\[
 \frac12\E[(\beta_k-\beta^*)^\top C(\beta_k-\beta^*)]
 =\frac12\operatorname{tr}(HU_k)
 =\frac12m_k^\top Hm_k+\frac12\operatorname{tr}(H\Sigma_k).
\]
These formulas require no stability assumption or inversion of
$I-\mathcal T_B$. For example, column vectorization gives the finite matrix
representation
\[
 \operatorname{vec}\{\mathcal T_B(U)\}=K_B\operatorname{vec}(U),\qquad
 K_B=R\otimes R+\frac{\eta^2}{B}
       \{\E[(zz^\top)\otimes(zz^\top)]-H\otimes H\},
\]
so \eqref{eq:discrete-second} is also an ordinary matrix-power and finite-sum
formula. Mean convergence alone does not ensure bounded second moments.

\subsection{A fixed-update rare-feature mean limit}

The common-block population setup also gives a finite-$k$ limit for both
GD and the mean of fresh-batch SGD. Index populations by $n$, not by the
number of observed training samples, and suppose
\[
 C_n=\begin{pmatrix}A_n&c_n\\c_n^\top&b_n\end{pmatrix}\succ0,\quad
 A_n\to A\succ0,\quad b_n=q_n(1-q_n)\to0,\quad q_n\to0,
 \quad c_n/b_n\to\kappa.
\]
The rare feature is a centered Bernoulli variable. Let
$S_n=\diag(S_{C,n},s_{R,n})$, where $S_{C,n}\to D$ is positive diagonal,
and suppose $b_n/s_{R,n}^2\to\rho\in[0,\infty)$.
Hold $\eta>0$ and $\beta^*$ fixed and take deterministic
$\beta_{n,0}\to\beta^0$ on the original scale. Set
\[
 G=D^{-2}A,\qquad F=I-\eta G,\qquad r_\eta=1-\eta\rho,
 \qquad d=\beta^*-\beta^0.
\]
\begin{proposition}[Rare-feature mean at fixed $k$]\label{prop:discrete-rare}
For each fixed integer $k\ge0$, the original-scale mean coefficients obey
\begin{align}
 \E\beta_{n,C,k}&\longrightarrow\beta_C^*-F^kd_C,\label{eq:discrete-rare-common}\\
 \E\beta_{n,R,k}&\longrightarrow\beta_R^*-r_\eta^kd_R
  +\eta\rho\sum_{j=0}^{k-1}r_\eta^{k-1-j}\kappa^\top F^jd_C.
  \label{eq:discrete-rare-coefficient}
\end{align}
For GD the expectation can be omitted. For SGD the assumptions of
Proposition~\ref{prop:discrete-moments} hold separately in each population;
any finite positive integer batch size $B_n$ gives the same mean.
\end{proposition}
\begin{proof}
The original-scale generator converges to a block triangular matrix:
\[
 P_nC_n\longrightarrow
 L=\begin{pmatrix}G&0\\\rho\kappa^\top&\rho\end{pmatrix}.
\]
For fixed $k$, matrix powering is continuous, so the mean error converges
to $(I-\eta L)^k(\beta^0-\beta^*)$.
The common block follows $F^k$, while solving the rare-coordinate scalar
recurrence driven by the common block gives
\eqref{eq:discrete-rare-coefficient}.
\end{proof}

The term $F^jd_C$ is the remaining common-coefficient error with its sign
reversed. Its convolution with $r_\eta^{k-1-j}$ transfers dependence through
$\kappa$ into the rare mean. For one common feature and standard-deviation
scaling, $G=1$ and $\rho=1$. From zero initialization, for $k\ge1$,
\[
 \E\beta_{n,R,k}\longrightarrow
   \{1-(1-\eta)^k\}\beta_R^*
   +\eta\kappa k(1-\eta)^{k-1}\beta_C^*.
\]
At $\eta=1,k=1$, the zeroth power is one; at $k=0$, the convolution is
zero. For fixed $k$, the general formula also covers $\rho=0$ (including
no rare-feature scaling), when the limiting rare coefficient stays at its
initial value. Variance scaling of the rare coordinate has
$b_n/s_{R,n}^2=1/b_n\to\infty$ and is outside this limit.

Binary feasibility is a distributional requirement, beyond specifying a
positive definite covariance. For common binary vector $W$ and rare
indicator $J$, one can choose distributions of $W\mid J=0$ and $W\mid J=1$
on $\{0,1\}^{p-1}$ and mix them with probabilities $1-q_n$ and $q_n$.
Then $c_n/b_n=\E(W\mid J=1)-\E(W\mid J=0)$ exactly, and nonnegative
conditional cell probabilities ensure feasibility. The numerical checks use
such supports, enumerate independent batches and short update trajectories,
and compare original- and scaled-coordinate moments with the formulas above.

The limit is only for each fixed update count (and hence any fixed finite
set of update counts). It is not uniform over unbounded $k$, does not pass
to an increasing stopping time or a fixed-data sampling theorem, and does
not assert convergence of second moments as $q_n\to0$. In particular,
fourth update moments and noise moments may grow with rarity despite a
finite limiting mean. Taking $\eta\to0$ with $k\eta\to T$ in the limiting
matrix power recovers the block population flow, a separate time-discretization
limit that supplies no joint guarantee for changing $n$, $k$, and $B_n$.
